\section{结论}
\label{sec:transferability}

$\mathrm{Ba_5Y_{12}Zn[O(SiO_4)]_8}$已有可靠的单晶结构和高温溶液合成证据，但尚无公开证据证明其批量相纯范围已被独立复现\cite{ababaikeri2024ba5y12zn}。因此，现阶段不宜给出“最佳配方”，应明确必须同时控制的变量和表征步骤。

相近体系给出三项一致认识。第一，Ba--Y--Si--O成相对组成坐标、MoO$_3$用量和冷却路径敏感；无Zn四方相还涉及非整比、竞争相和玻璃区\cite{wierzbickawieczorek2017high,yamane2024synthesis,gulay2024navigation}。第二，陶瓷样品中的副相会改变名义组成的解释，单晶身份不能代替批量相定量\cite{yamane2024microstructure}。第三，Ba--Zn--Si--O的晶型和相稳定性随热史和取代组成变化，因此Zn源、冷却路径和观察温度必须联合记录\cite{lin1999phase,kerstan2013bazn2si2o7,zou2019anti}。这些认识可指导实验设计，但不能把近邻数值直接转写成目标相条件。

下一阶段最有信息量的工作依次是：(i)以BaCO$_3$、Y$_2$O$_3$、ZnO和SiO$_2$为基础原料，在目标组成附近建立局部相图；(ii)检验$\mathrm{Ba_5Y_{13-x}Zn_x[SiO_4]_8O_{8.5-x/2}}$系列，区分有限固溶、线化合物和两相区；(iii)并行开展固相成相和高温溶液长晶，使PXRD/Rietveld回答批量相区，SCXRD回答结构身份；(iv)目标相获得独立重复后，再用Mg/Co类比检验尺寸差异。逆合成模型推荐的氧化物或碳酸盐组合可作为前驱体起点，但不能替代相图、热史和结构验证。

最值得立即开展的下一步实验是目标组成附近的小步长局部相图：以BaCO$_3$、Y$_2$O$_3$、ZnO和SiO$_2$为前驱体，在相近体系的温区内分段煅烧和退火，并对每个组成点进行定量PXRD/Rietveld、EDS或EPMA，优先确定目标相是否存在可重复的批量相区。随后将近单相组成转入高温溶液长晶，用SCXRD复核结构。这样的实验在成功或失败时都能更新假设。小步长组成网格、实际元素和氧计量、单晶与粉末衍射的交叉验证以及完整保存的负结果，可形成可迭代的实验依据，并回答Zn是否进入无Zn四方谱系、目标相是否具有可重复的稳定区，以及晶体生长条件与批量成相条件是否属于同一相平衡区域。
